\chapter{Rastreabilidade da implementação e das evidências}
\label{ap:rastreabilidade}

A campanha principal é identificada como E11-M01 no acervo; seu pré-registro é E11-M00-v1, e o inventário inicial do ambiente é E11-D01. Este apêndice reúne os caminhos da implementação completa, preservada na branch \texttt{master} do \href{https://github.com/JnetSilvestre/tcc1_JoaoVictor_BigData/tree/13d1892da1be68203e0693beb19e2b2c11e7c512}{repositório do trabalho, na versão de integração da monografia}. A branch \texttt{tcc2/overleaf-utfpr} contém os arquivos necessários à compilação; nela, as cópias dos códigos citados estão em \path{uteis/codigo-original/} e \path{uteis/codigo-tcc2/}. Os demais caminhos da tabela referem-se à implementação completa. A Tabela~\ref{tab:mapa-evidencias} permite localizar o material sem depender de um caminho absoluto do computador utilizado.

\begin{table}[htbp]
\centering\small
\caption{Mapa entre afirmações e artefatos de suporte.}
\label{tab:mapa-evidencias}
\begin{tabularx}{\linewidth}{@{}p{4.1cm}X@{}}
\toprule
Afirmação ou procedimento & Artefato na implementação completa \\
\midrule
Origem do núcleo Python & \path{docs/tcc2/codigo-original/} e sua nota de proveniência \\
Adaptação de entidades & \path{prova-conceito/pon/} \\
Referência local e oráculo & \path{prova-conceito/c0/pon_execution.py} e \path{prova-conceito/c0/oracle.py} \\
Estado e contratos HTTP & \path{prova-conceito/c1/state.py} e \path{prova-conceito/c1/protocol.py} \\
Imagem e topologia kind & \path{prova-conceito/Dockerfile} e \path{infraestrutura/kind/kind-config.yaml} \\
Implantação funcional C3/C4 & \path{infraestrutura/kubernetes/c3/manifests.yaml} e \path{infraestrutura/kubernetes/c4/manifests.yaml} \\
Agenda e ciclo da campanha & \path{prova-conceito/benchmark/m01_campaign.py} e \path{prova-conceito/benchmark/m01_lifecycle.py} \\
Auditoria e extração & \path{prova-conceito/benchmark/m01_extract.py} \\
Métricas, tabelas e gráficos & \path{prova-conceito/benchmark/m01_analyze.py} e \path{prova-conceito/benchmark/m01_publication.py} \\
Laudo da amostra & \path{docs/tcc2/etapa-12-auditoria-extracao-e11-m01.md} \\
\bottomrule
\end{tabularx}
\fonte{}
\end{table}

O código de referência foi fixado na revisão
\begin{center}\small\texttt{c2d525b5c675820806af34348e259bff32a6bc5f}.\end{center}
A revisão do controlador da série é
\begin{center}\small\texttt{24740e922181dace018a26938b353f10d952da06}.\end{center}
A rodada B01/C2 conserva a revisão de execução
\begin{center}\small\texttt{eecc15e3bb1c36c9ee51b3165f882d76a0c6ed9b}.\end{center}
Esses identificadores representam papéis distintos: origem do núcleo, controlador da campanha e execução anterior da primeira rodada. A revisão editorial da monografia não substitui os identificadores registrados nas evidências.

O pacote \path{E11-M01-auditoria-extracao.zip} contém os arquivos derivados necessários à análise. Seu SHA-256 é apresentado em duas linhas apenas para leitura; a verificação utiliza a sequência contínua de 64 caracteres:
\begin{center}\ttfamily\small
66d115c77f7b22819cc43ad84e11ccd43\\
759ca319ac5ef26c3f23e2af8645329
\end{center}

Os ZIPs originais e os journals permanecem no acervo de evidências. O pacote derivado permite reproduzir as estatísticas publicadas, enquanto a repetição integral da auditoria de logs, estado e topologia requer esse acervo original. Os 30 registros históricos de revisão do controlador são distinguidos da auditoria binária posterior no laudo da Etapa 12.

\section*{Exemplo de topologia observada em uma rodada real}

A Tabela~\ref{tab:topologia-real} transcreve campos do arquivo \texttt{topology-end.json} arquivado na rodada C4 do bloco B10. O recorte foi extraído do ZIP original e seu hash interno foi conferido. Trata-se de evidência histórica da rodada, e não de uma consulta atual ao cluster ou de uma representação de saída simulada.

\begin{table}[htbp]
\centering\small
\caption{Posicionamento registrado ao final da rodada C4/B10.}
\label{tab:topologia-real}
\begin{tabularx}{\linewidth}{@{}p{3.8cm}Xr@{}}
\toprule
Papel & Nó registrado & Reinícios \\
\midrule
\texttt{attribute-source} & \texttt{tcc2-pon-worker} & 0 \\
\texttt{logic-engine} & \texttt{tcc2-pon-control-plane} & 0 \\
\texttt{rule-executor} & \texttt{tcc2-pon-worker} & 0 \\
\bottomrule
\end{tabularx}
\fonte{Arquivo original \texttt{topology-end.json}, rodada C4/B10}
\end{table}

Os nomes e UIDs completos dos Pods, os identificadores de contêiner e os demais campos permanecem no arquivo documental \path{docs/tcc2/e11-m01-exemplo-c4-b10-topology-end.json}. O inventário do pacote derivado liga essa rodada ao ZIP original de SHA-256:
\begin{center}\ttfamily\small
a21bd54e385d70923ed2e6b84e7be6a46\\
a41a592fdd5257cdf5385927978b24d
\end{center}
Essa evidência exemplifica a diferença entre declarar um seletor de nó no manifesto e verificar o posicionamento efetivo dos componentes. A conclusão sobre a campanha completa depende da auditoria das 30 rodadas, e não apenas deste exemplo.
